\section{Introduction}
Many knowledge-based and perception systems represent their input as a grid of vectors: the embeddings of the patches of an image, the feature vectors of a sensor array, or the orientation field of an ordered material. Unsupervised anomaly detection on such data usually asks how unusual the vector of a site is, for example by its distance to the nearest members of a memory of normal embeddings \citep{Roth2022,Defard2021} or by density- and isolation-based outlier scores \citep{Breunig2000,Liu2008}; surveys cover the many variants \citep{Chandola2009,Ruff2021,Pang2021,Aggarwal2017}. Such scores are blind by construction to an anomaly that consists only of ordinary vectors in the wrong places. Detectors for \emph{spatial} or \emph{contextual} outliers compare a site with its neighbors instead \citep{Shekhar2003,Song2007,Schubert2014}, and benchmarks of logical anomalies in images show that arrangement errors are a practically relevant class \citep{Bergmann2022}. The two families are usually developed separately, and how they should be combined, and when either can work at all, is rarely made explicit.

This paper studies that question on lattices of embeddings with a combination of elementary analysis and controlled computation. It addresses three research questions.
\begin{description}
\item[RQ1] Under what conditions does a score that depends only on a site's embedding fail to detect an anomaly that a spatial-order statistic can detect, and vice versa?
\item[RQ2] How do the geometry of an anomaly, the size of the neighborhood and the coherence of the background determine whether spatial aggregation preserves an anomaly and whether a spatial-order statistic can see it?
\item[RQ3] Can a label-free, site-by-site fusion of the two kinds of score detect feature and arrangement anomalies that occur \emph{in the same lattice} more reliably than choosing one score per cohort, and is a locally adaptive fusion worth its complexity?
\end{description}

The study builds on an earlier version of this work, which proposed a similarity-weighted rarity score (channel A), a directional order statistic (channel B) and a cohort-level choice between them by a tail index, and evaluated them on scenarios of which most had been used during development. That design was criticized, rightly, for the dependence between development and test data, for the lack of independent validation and because a cohort-level choice cannot serve lattices that contain both kinds of anomaly. The present version keeps the earlier results only as motivation and replaces the evaluation.

\textbf{Contributions.} (i) Five analytical statements (Section~\ref{sec:theory}): position-free scores are blind to rearrangements; a fixed aggregation window preserves an anomalous site only if its local anomaly occupancy exceeds a threshold computable from the scores; the directional order statistic is insensitive to straight walls between ordered domains; a site-wise maximum of calibrated scores retains every detection of either channel at the cost of at most their combined false positives; and rotation-invariant local statistics cannot see the interior of a rotated patch. (ii) A detectability map over feature strength, arrangement disruption and background coherence (Experiment A). (iii) A comparison of seven label-free fusion rules on lattices with both kinds of anomaly, including two locally adaptive rules (Experiment B). (iv) A two-stage pre-registered evaluation on generators that were frozen before they were scored (Experiment C): the first frozen rule failed, for a reason that we identify, and a revised rule was then confirmed on a third generator. (v) Robustness and runtime measurements (Experiment D). The results are reported whether or not they favor the methods proposed in the earlier version, and several do not. All data are synthetic; Section~\ref{sec:discussion} states what this does and does not establish.

\section{Related work}
\textbf{Feature-space anomaly scores.} Nearest-neighbor distances to a memory bank of normal patch embeddings are the basis of strong industrial inspection methods \citep{Roth2022}; Gaussian models of patch embeddings \citep{Defard2021}, LOF \citep{Breunig2000} and isolation forests \citep{Liu2008} are other marginal scores. All of them assign a score to a vector, not to its placement, which is the property formalized in Proposition~\ref{prop:perm}.

\textbf{Spatial, contextual and logical anomalies.} Spatial outlier detection compares an attribute with the aggregate of its spatial neighbors \citep{Shekhar2003}; conditional (contextual) anomaly detection models behavioral attributes given contextual ones \citep{Song2007}; Schubert et al.~\citep{Schubert2014} unify local outlier methods by their notion of locality. In vision, logical anomalies, such as missing, misplaced or wrongly combined parts, have their own benchmark \citep{Bergmann2022} and dedicated models \citep{Batzner2024}. The directional statistic studied here is a contextual score of a specific kind: it measures the dispersion of the unit vectors around a site \citep{Mardia2000,Banerjee2005} and takes the minimum over four sub-windows, as the Kuwahara filter does for intensity variance \citep{Kuwahara1976}, so that walls between ordered domains do not count as disorder.

\textbf{Score aggregation and fusion.} Smoothing per-site scores over a window is routine in anomaly localization; edge-preserving weights \citep{Tomasi1998} avoid averaging across boundaries. Combining detectors is the subject of outlier ensembles \citep{Aggarwal2015}, where maximum and average combination are standard, and comparing scores from different detectors requires their calibration to a common scale \citep{Kriegel2011}; here cohort ranks provide that scale without labels. The contribution of this paper is not a new ensemble rule but the analysis of when the two channels are complementary, when aggregation destroys a channel's signal, and what a site-wise maximum costs.
